\chapter{Stability of Matter from Hardy}\label{ch:som}
\index{stability of matter}\index{Hardy inequality}\index{Lieb--Thirring inequality}\index{Dyson--Lenard proof}\index{Thomas--Fermi scaling}\index{free-boundary problem}\index{Bernoulli condition}\index{Alt--Caffarelli theorem}\index{Caffarelli--Lin theory}\index{viscosity solutions}\index{level-set method}\index{computer-assisted proof}\index{form-bounded potential}\index{Banach--Alaoglu theorem}\index{Rellich--Kondrachov theorem}\index{concentration--compactness}

\begin{quote}
\itshape I hope later to make clear that the reigning doctrine is born of distress.\\
\normalfont\hfill --- E.~Schr\"odinger, \emph{The Present Situation in Quantum Mechanics}, 1935.
\end{quote}
\bigskip

Stability of matter --- the requirement that the total electronic ground-state energy of a system of $N$ atoms be bounded below by an extensive quantity, linear in $N$ rather than diverging more rapidly --- is among the most fundamental properties demanded of any quantum theory of matter. Without it, bulk matter could not exist. A theory that allows the energy of a hundred atoms to be ten times more negative than the energy of ten atoms is internally consistent; a theory that allows the energy of a hundred atoms to be a hundred-fold more negative than ten atoms predicts collapse. Matter exists because the latter does not happen.

Within standard quantum mechanics, proving stability of matter has been a major mathematical undertaking spanning four decades and several hundred pages. Within the multiphase real-space formulation, stability follows directly from Hardy's inequality in approximately one paragraph. The contrast is the subject of this chapter.

We argue that the proof difficulty in StdQM is not a feature of the underlying physics but a symptom of the formulation: working in $3N$-dimensional configuration space with overlapping orbitals forces the singularities of the Coulomb potential to be controlled globally, requiring elaborate apparatus. Working in real 3D space with non-overlapping unit densities localises the control, and stability falls out of the variational structure.

\section{The classical problem: Dyson--Lenard--Lieb--Thirring}\label{sec:som-classical}

The question \emph{is matter stable?} is older than quantum mechanics itself. The classical instability is straightforward: a system of point charges interacting through Coulomb forces collapses, because the potential energy diverges as $-1/r$ when two opposite charges approach. Classically, only the uncertainty principle prevents this collapse; the kinetic-energy cost of localising an electron to a small region balances the potential-energy gain.

In quantum mechanics, the question becomes \emph{is the ground-state energy of a many-electron Coulomb system bounded below by an extensive quantity?} The naive answer is yes (the uncertainty principle is operative quantum-mechanically) but the proof for an $N$-electron system is not trivial. It requires showing that as the number of electrons grows, the total energy grows linearly in $N$, not as a higher power.

\paragraph{Dyson--Lenard (1967).}
\textbf{Dyson and Lenard} gave the first rigorous proof in 1966~\cite{DysonLenard67}, a 26-page paper that Dyson himself characterised as ``awful mathematics''. The proof established stability of matter for fermionic Coulomb systems, but the bound it produced --- the constant in $E \ge -CN$ --- was astronomically larger than the physical scale.

\paragraph{Lieb--Thirring (1975).}
\textbf{Lieb and Thirring} compressed the proof to three pages in 1975~\cite{LiebThirring75}, using the kinetic-energy inequality that bears their name. The Lieb--Thirring bound on the ground-state energy of an atom of nuclear charge $Z$ scales as
\begin{equation}\label{eq:lieb-thirring}
E(Z) \;\sim\; -Z^{7/3},
\end{equation}
a result derived from the Thomas--Fermi model rather than the full Schr\"odinger equation. The proof and its extensions occupy the 300-page book \emph{The Stability of Matter in Quantum Mechanics}~\cite{LiebSeiringer05}.

\paragraph{Where the difficulty comes from.}
The proof's central difficulty stems from the \textbf{global support of electron wave functions} in StdQM. In principle, every electron can be everywhere; every electron can interact with every nucleus and with every other electron. Intricate book-keeping is required to prevent local accumulation of charge density from driving the potential energy to $-\infty$. The wave function on $\mathbb{R}^{3N}$ has to be controlled globally, and the kinetic-energy term has to be shown to overwhelm the worst-case potential-energy accumulation in every region simultaneously.

The fermionic antisymmetry of the wave function is essential to the proof: it is what prevents an arbitrary number of electrons from being concentrated in a small region. The Lieb--Thirring inequality is a quantitative form of this statement, bounding the kinetic energy of an antisymmetric wave function from below by an integral of the one-body density to a fractional power.

\paragraph{Status of the proof.}
The proof is not part of standard textbook treatments of quantum mechanics, even though stability is a precondition for everything those textbooks describe. A student of standard QM learns to compute atomic energies and chemical bonds without learning that the underlying theory has been shown not to collapse; the proof is a specialist topic in mathematical physics.

\section{Hardy's inequality}\label{sec:hardy}

The multiphase formulation reaches stability of matter through a different inequality. For any $\psi \in H^1(\mathbb{R}^3)$, the \textbf{Hardy inequality} states
\begin{equation}\label{eq:hardy}
\int_{\mathbb{R}^3} \frac{\psi^2(\mathbf{x})}{|\mathbf{x}|^2}\, d\mathbf{x}
\;\le\; 4 \int_{\mathbb{R}^3} |\nabla \psi(\mathbf{x})|^2\, d\mathbf{x}.
\end{equation}
A related form, more directly relevant here, bounds the Coulomb-singular potential energy by a product of the $L^2$-norm and the kinetic-energy semi-norm:
\begin{equation}\label{eq:hardy-kato}
\int_{\mathbb{R}^3} \frac{\psi^2(\mathbf{x})}{|\mathbf{x}|}\, d\mathbf{x}
\;\le\;
\Bigl(\int_{\mathbb{R}^3} \psi^2\, d\mathbf{x}\Bigr)^{1/2}
\Bigl(\int_{\mathbb{R}^3} |\nabla \psi|^2\, d\mathbf{x}\Bigr)^{1/2}.
\end{equation}
This bound was used by Kato in his analysis of the Schr\"odinger equation; it expresses the same regularity content as the standard Hardy inequality~\eqref{eq:hardy} and is more conveniently applied to the Coulomb potential.

The Hardy inequality is a regularity statement: a function in $H^1(\mathbb{R}^3)$ cannot concentrate its $L^2$-mass too close to the origin, because doing so would force its $H^1$ semi-norm to blow up. The Coulomb singularity at the origin is integrable against an $H^1$ density precisely because the kinetic-energy semi-norm controls the concentration.

\section{The one-paragraph derivation}\label{sec:som-derivation}

We now derive stability of matter for RealQM. The argument has three steps.

\paragraph{Step 1: atomic bound.}
In RealQM, the total wave function $\Psi$ of an atom with nuclear charge $Z$ has unit-density components meeting at the Bernoulli free boundary (continuity plus homogeneous Neumann), with $\Psi \in H^1(\mathbb{R}^3)$ and $\|\Psi\|_2^2 = Z$ (one unit-charge density per electron, $Z$ electrons in a neutral atom). The kernel--electron potential energy is
\[
PK \;=\; -Z \int \frac{\Psi^2(\mathbf{x})}{|\mathbf{x}|}\, d\mathbf{x}.
\]
Applying~\eqref{eq:hardy-kato} to the global atomic $\Psi$,
\[
\int \frac{\Psi^2(\mathbf{x})}{|\mathbf{x}|}\, d\mathbf{x} \;\le\; \sqrt{Z} \cdot \sqrt{K(\Psi)},
\]
where $K(\Psi) = \tfrac{1}{2}\int |\nabla\Psi|^2$ is the RealQM kinetic-energy term, up to a constant factor. The total energy is therefore bounded below by
\[
E_{\rm atom} \;\ge\; K - Z \sqrt{Z} \sqrt{K}.
\]
Minimising the right-hand side over $K > 0$ gives the optimal kinetic-energy choice $\sqrt{K} \propto Z^{3/2}$, and the atomic energy lower bound
\begin{equation}\label{eq:atom-bound}
E_{\rm atom} \;\ge\; -C\, Z^3,
\end{equation}
with $C$ an absolute constant.

\paragraph{Step 2: bulk bound.}
For a system of $N$ non-overlapping atoms each of nuclear charge $Z$, the electron territories of distinct atoms are non-overlapping in physical space by the multiphase formulation. The Hardy bound~\eqref{eq:hardy-kato} applied to each atom involves only the electronic density on that atom's territory. The bounds therefore add atom by atom:
\begin{equation}\label{eq:bulk-bound}
E \;\ge\; -C\, N\, Z^3.
\end{equation}
This is the ``stability of the second kind'' of the Dyson--Lenard--Lieb--Thirring tradition: total energy is extensive in $N$, with no possibility of collapse to $-\infty$.

\paragraph{Step 3: that is the proof.}
\textbf{The non-overlap constraint is the physical origin of stability in RealQM.} Each atom's potential-energy contribution is controlled by its own kinetic energy through Hardy on its own region of support; no global book-keeping is required because the contributions from distinct atoms simply add. Stability is a direct consequence of the formulation rather than a deep theorem to be proved separately.

That is the entire proof. Three steps: Hardy on each atom, addition over atoms, done.

\section{Why the RealQM proof is so much shorter}\label{sec:why-shorter}

The contrast with Dyson--Lenard--Lieb--Thirring is informative. The StdQM proof is hard because:
\begin{itemize}
\item Electron wave functions in StdQM have global support on $\mathbb{R}^{3N}$. Any electron can be anywhere; the potential singularities at every nucleus must be controlled simultaneously.
\item Antisymmetry is essential to the proof and has to be invoked through the Lieb--Thirring inequality, which is itself a non-trivial functional-analytic result.
\item The Thomas--Fermi--like $-Z^{7/3}$ scaling has to be reached through a careful minimisation argument that involves the entire many-body kinetic-energy operator on antisymmetric subspaces.
\end{itemize}
The RealQM proof is short because:
\begin{itemize}
\item Electron wave functions in RealQM have disjoint supports by construction. Each electron has its own territory; the potential singularities at distinct nuclei are automatically separated.
\item Antisymmetry is replaced by non-overlap, which is a geometric constraint applied at the level of supports rather than at the level of permutation symmetry. The geometric constraint is what makes the global bound separate into a sum of local bounds.
\item Hardy's inequality is a standard regularity result that applies to any $H^1$ function on $\mathbb{R}^3$; it does not require antisymmetry, does not require fermion statistics, and does not need a separate variational subspace.
\end{itemize}

The price for the shorter proof is the trade-off discussed in Chapter~\ref{ch:equation}: antisymmetry is replaced by non-overlap, with the two requirements not being equivalent. The bound~\eqref{eq:bulk-bound} is proved for a system of non-overlapping unit-density wave functions, and that is what RealQM describes. It is not proved for arbitrary antisymmetric wave functions on $\mathbb{R}^{3N}$, and the StdQM bound therefore remains the right result for the StdQM formulation.

\section{Atomic scaling: heuristic $-Z^2 \log Z$ within RealQM}\label{sec:atom-scaling}

The rigorous Hardy bound~\eqref{eq:atom-bound} establishes stability with $E_{\rm atom} \ge -CZ^3$, but the bound is not tight; the actual ground-state energy of an atom grows more slowly with $Z$. A heuristic estimate within the unit-density formulation can be obtained from an approximate shell structure.

\paragraph{Shell-scaling heuristic.}
Suppose the $m$-th filled shell of an atom has:
\begin{itemize}
\item Radius $r_m \sim m^3 / Z$ (Bohr-like scaling adjusted for shell index).
\item Thickness $d_m \sim m^2 / Z$.
\item Population $2m^2$ electrons in the $m$-th filled shell.
\item Charge density $\rho_m(r) \sim Z / r^2$ in spherical symmetry.
\end{itemize}
Integrating kinetic and potential contributions over a density profile of this form produces, at heuristic level, an atomic ground-state energy of order
\begin{equation}\label{eq:heuristic-scaling}
E(Z) \;\sim\; -\log(Z)\, Z^2,
\end{equation}
i.e.\ $-Z^2$ up to a logarithmic correction.

\paragraph{Status of the heuristic.}
The shell-scaling laws above are \emph{heuristic} estimates, not rigorously derived asymptotics, and the resulting $-Z^2 \log Z$ form is intended only as an order-of-magnitude indication. The Thomas-Fermi--based Lieb--Thirring scaling $-Z^{7/3}$ is a different heuristic prediction; both lie comfortably below the rigorous $-Z^3$ Hardy upper bound on $|E|$ and disagree only at the level of subleading corrections.

A precise comparison of the two heuristic scalings, and of either with experimental atom energies across the periodic table, is left for follow-up work. What the heuristic does establish is that the Hardy bound is not asymptotically tight: the actual atomic energies grow more slowly than $Z^3$, and the slower growth is consistent with the Level-1 numerical results of Chapter~\ref{ch:atoms} across Li--Rn.

\section{Kinetic-energy upper bound and existence of a minimiser}\label{sec:existence}

The Hardy bound of Section~\ref{sec:som-derivation} controls the total energy from below. The same inequality also controls the kinetic energy from \emph{above} at any bound state, and the combination is what delivers existence of a minimiser through a standard compactness argument.

\paragraph{Upper bound on the kinetic energy from Hardy.}
The argument is short. The multiphase wave function $\Psi = \sum_i \psi_i$ has $\|\Psi\|_2^2 = N$ (each $\psi_i$ unit-normalised on $\Omega_i$, supports disjoint, so cross-terms vanish), and $\|\nabla \Psi\|_2^2 = 2K$. Applying~\eqref{eq:hardy-kato} to $\Psi$ at each kernel position $\mathbf{R}_a$ and summing:
\begin{equation}\label{eq:Vbound}
|V_{eK}| \;=\; \sum_a Z_a \int \frac{\Psi^2}{|\mathbf{x} - \mathbf{R}_a|}\, d\mathbf{x}
\;\le\; \sqrt{N}\,\Bigl(\textstyle\sum_a Z_a\Bigr)\, \sqrt{2K}.
\end{equation}
The inter-electron repulsion is non-negative ($V_{ee} \ge 0$). Combining with the energy decomposition $E = K + V_{eK} + V_{ee}$,
\begin{equation}
K \;=\; E - V_{eK} - V_{ee} \;\le\; E + \sqrt{N}\,\Bigl(\textstyle\sum_a Z_a\Bigr)\sqrt{2K}.
\end{equation}
Setting $u = \sqrt{2K}$ and writing $C_N = \sqrt{N}\sum_a Z_a$ for compactness, the inequality reads $\tfrac{1}{2}u^2 - C_N\, u - E \le 0$, which for any finite $E$ implies
\begin{equation}\label{eq:Kbound}
K \;\le\; \tfrac{1}{2}\Bigl(C_N + \sqrt{C_N^2 + 2E}\Bigr)^{\!2}, \qquad C_N \equiv \sqrt{N}\,\textstyle\sum_a Z_a.
\end{equation}

The structural relation behind this bound is that the Hardy inequality expresses the Coulomb potential as \textbf{form-bounded by the kinetic energy with sub-linear factor}: $|V| \le c\sqrt{K}$. This is stronger than the relative-form-bound condition $|V| \le c K$ with $c < 1$ used in Kato--Rellich self-adjointness; the sub-linear scaling means $|V|/K \to 0$ as $K \to \infty$, which is the form-boundedness with \textbf{zero relative bound} familiar from the spectral theory of Schr\"odinger operators. The kinetic energy therefore cannot run away to $+\infty$ while the total energy $E = K + V$ remains finite: any growth of $K$ is matched at most by a $\sqrt{K}$-scale decrease in $V$, so the sum keeps a finite minimum.

For the H atom ($N = 1$, $Z = 1$, $E = -\tfrac{1}{2}$), one has $C_N = 1$ and the bound~\eqref{eq:Kbound} gives $K \le \tfrac{1}{2}$ exactly, saturated by the actual ground-state value $K = \tfrac{1}{2}$. For multi-electron atoms and molecules the bound is uniform but loose; it is independent of the variational status of the wave function, holding equally at constrained minima where the virial relation $2K + V = 0$ does not hold.

\paragraph{Existence of a ground state.}
The kinetic-energy bound~\eqref{eq:Kbound} guarantees that any energy-minimising sequence is bounded in $H^1(\mathbb{R}^3)$. The $L^2$ norm is fixed by the unit-charge constraint, $\|\Psi_n\|_2^2 = N$ for every $n$, and the kinetic semi-norm $\|\nabla \Psi_n\|_2 = \sqrt{2K(\Psi_n)}$ is bounded uniformly along the sequence because the energy stays bounded above (by definition of a minimising sequence). Banach--Alaoglu then yields a subsequence $\Psi_n \rightharpoonup \Psi^*$ converging weakly in $H^1(\mathbb{R}^3)$, and weak lower semi-continuity of $K$ together with strong $L^2$-convergence on bounded sets (Rellich--Kondrachov) gives $E(\Psi^*) \le \liminf E(\Psi_n) = \inf E$. So a candidate minimiser $\Psi^*$ exists in the weak sense.

A complete existence proof for the full multiphase variational problem requires two additional analytical steps not delivered by the H-bound alone: \textbf{(i)} controlling possible loss of mass at infinity along the sequence (handled by concentration--compactness arguments of Lions type for atomic and molecular systems), and \textbf{(ii)} verifying that the weak limit $\Psi^*$ inherits the multiphase structure of disjoint supports, unit-norm per electron, and the Bernoulli boundary conditions of Chapter~\ref{ch:equation}, Section~\ref{sec:bernoulli}. Both are technically non-trivial but use standard tools from the calculus of variations and free-boundary theory, and are left here to follow-up analytical work; the kinetic-energy bound from Hardy is the load-bearing first ingredient.

\section{Connection to the free-boundary literature}\label{sec:freebdy-lit}

The Bernoulli condition at the inter-electron boundary $\Gamma_i$ places the multiphase variational problem in a \textbf{well-studied class of free-boundary problems}, with an established analytical literature on existence and regularity. We mention four results that bear most directly on RealQM.

\begin{itemize}
\item The foundational \textbf{one-phase Bernoulli problem} of Alt and Caffarelli~\cite{AltCaffarelli81}: minimisers of $\int |\nabla u|^2 + \lambda^2 \mathbf{1}_{\{u>0\}}$ over admissible $u$ exist and are Lipschitz, with the free boundary $\partial\{u > 0\}$ analytic outside a small singular set. The structure --- $H^1$ minimiser, Lipschitz regularity inside the support, analytic free boundary --- is exactly what RealQM produces for the one-electron pieces $\psi_i$ on $\Omega_i$.

\item The \textbf{two-phase Bernoulli problem} of Alt, Caffarelli, and Friedman~\cite{AltCaffarelliFriedman84}: when the admissible functions are allowed to change sign with one phase on each side of the free boundary, the analogous regularity holds with a more delicate analysis of the interface. This is the structurally closest classical problem to a pair of touching RealQM electron territories.

\item The \textbf{multi-phase competing-system theory} of Caffarelli and Lin~\cite{CaffarelliLin08}: a finite collection of non-negative functions $\{u_i\}$ minimising $\sum_i \int |\nabla u_i|^2 + F$ subject to a disjoint-supports constraint $u_i u_j = 0$ for $i \ne j$ is shown to have $C^{0,1}$ regular components and a free-boundary structure with singular set of low Hausdorff dimension. This is the closest classical analogue of the full RealQM multiphase setup: the disjoint-supports constraint matches the non-overlap of RealQM electron territories.

\item The \textbf{spectral minimal partition} theory of Helffer, Hoffmann-Ostenhof, and Terracini~\cite{HHOT09}: minimising the sum of Dirichlet eigenvalues over partitions of a domain into $N$ regions produces a partition with the same Bernoulli-type interface structure. This frames the RealQM partition $\{\Omega_i\}$ as a coupled-Coulomb generalisation of a spectral minimal partition.
\end{itemize}

The shared structural feature across these problems is that the partition or support pattern is the \emph{primary} degree of freedom, varying jointly with the function(s) supported on it, and the boundary conditions at the free interface emerge from the variational principle rather than being externally imposed. RealQM differs from the autonomous multiphase systems of Caffarelli--Lin in one structural respect: the coupling between phases via the Hartree potentials $V_j$ in the per-phase Hamiltonian~\eqref{eq:Hi} of Chapter~\ref{ch:equation} is non-local in the bulk variables. Adapting the regularity results above to this coupling is the natural analytical route to a complete regularity theory for $\Gamma_i$; the existing results provide the framework and the techniques.

\section{Computational versus analytical existence and regularity}\label{sec:comp-vs-anal}

The analytical results assembled above --- the Hardy stability bound~\eqref{eq:bulk-bound}, the boundedness of minimising sequences in $H^1(\mathbb{R}^3)$ from the kinetic-energy bound~\eqref{eq:Kbound}, and the free-boundary regularity transferable from~\cite{AltCaffarelli81, AltCaffarelliFriedman84, CaffarelliLin08, HHOT09} --- establish the analytical scaffolding for existence and regularity of the RealQM ground state in a weak sense. What remains for a complete analytical theory is the adaptation of these results to the specific RealQM setting: Coulomb singularities at each nucleus (handled locally by Hardy), and the coupling between phases through the Hartree potentials in~\eqref{eq:Hi}, which is the structural feature distinguishing RealQM from the autonomous multiphase systems of Caffarelli--Lin. These adaptations are routine in scope but technically substantial, and have not been written down in full anywhere.

\paragraph{Level-set route.}
A complementary, and in practice more direct, route to existence and regularity is the \textbf{level-set formulation} of the relaxation dynamics (Chapter~\ref{ch:relaxation}, Section~\ref{sec:three-line}). The boundary $\Gamma_i$ is encoded implicitly as the zero level set of an auxiliary function $\chi_i$ on $\mathbb{R}^3$, evolved by the Hamilton--Jacobi-type equation $\dot\chi_i + \beta_i |\nabla \chi_i| = 0$ from~\eqref{eq:3line-chi}. Smoothness of $\Gamma_i$ is then \emph{inherited} from smoothness of $\chi_i$ via the implicit function theorem: where $\nabla \chi_i \ne 0$, the level set $\{\chi_i = 0\}$ is automatically a smooth manifold; singular points of $\Gamma_i$ correspond exactly to $\nabla \chi_i = 0$ and are detectable directly in the simulation. The level-set framework therefore reformulates the regularity question, shifting it from ``is $\Gamma_i$ smooth?'' to ``is $\chi_i$ smooth?'', and converts the abstract regularity question into a constructive observable on the simulation grid. The analytical theory of level-set evolution equations (\textbf{viscosity solutions} in the sense of Crandall--Lions and Evans--Spruck~\cite{CrandallLions83, EvansSpruck91}) provides the rigour for the level-set side; the simulation provides explicit smoothness for any specific configuration. This is methodologically distinct from but operationally equivalent to the regularity programme of~\cite{AltCaffarelli81, CaffarelliLin08}, and for RealQM's practical purposes it is the more direct route.

\paragraph{Qualitative analytical versus quantitative computational.}
Beyond the level-set formulation, analytical theorems on free-boundary regularity typically deliver \emph{qualitative} statements --- the free boundary is $C^{1,\alpha}$ off a singular set of Hausdorff dimension at most $n-2$, with exponent $\alpha$ unspecified, the constants astronomical, the singular-set location not located. A converged simulation by contrast delivers, for the specific configuration, the exact locations of triple-junction and quadruple points, explicit H\"older exponents at smooth boundary points, quantitative Lipschitz constants for the normal field along $\Gamma_i$, density ratios at boundary points, and direct verification that $|\partial_n \psi_i| \to 0$ as the simulation converges. The three-line code of Section~\ref{sec:three-line}, applied to atoms, molecules, condensed phases, and protein-scale systems, reaches converged solutions in which the wave functions $\psi_i$ stabilise to spatial profiles consistent with~\eqref{eq:Hi}--\eqref{eq:neumann}, the Bernoulli free boundary stabilises with continuity and homogeneous Neumann holding to mesh resolution, and the total energy plateaus. \textbf{Each such converged simulation may be regarded as a constructive proof, to a specified numerical precision, that a solution of the variational problem exists, that the wave functions are Lipschitz, and that the free boundary is regular off a small singular set whose location is explicitly determined.} The \emph{quantitative} content of this constructive evidence routinely exceeds the qualitative content of the analytical theorems: a chemist who needs to know the angle at which $\Gamma_i$ meets a bond axis can read it off the simulation but cannot extract it from any qualitative regularity theorem.

\paragraph{The two routes are complementary.}
The analytical scaffolding gives the universal qualitative picture: stability holds, minimisers exist in a weak sense, free boundaries are regular off small singular sets, and the underlying mechanism is encoded in the Bernoulli structure of Chapter~\ref{ch:equation}, Section~\ref{sec:bernoulli}. The computational diagnostics give the specific quantitative picture for each system: where the boundaries sit, where the singular points live, with which numerical regularity exponents. Neither is complete on its own; together they cover both the principled and the operational regimes. This division of labour is in the spirit of \textbf{computer-assisted proofs} (the 4-colour theorem, the Kepler conjecture, parts of mathematical fluid dynamics), in which finite computational verification, systematically combined with analytical reduction, replaces traditional proof in cases where the latter is intractable for the geometry at hand. The Gallery's runnable scan files are the operational embodiment of this hybrid framework: anyone doubting the regularity properties of a specific RealQM solution can run the corresponding file and inspect the boundary directly, with the analytical scaffolding established in the cited literature ensuring that what they see is what the variational principle predicts.

\section{Significance: stability as an introductory result}\label{sec:som-significance}

Stability of matter, viewed in StdQM as one of the deepest theorems of mathematical physics, becomes in RealQM a one-paragraph consequence of Hardy's inequality applied to non-overlapping electronic territories. The contrast is informative beyond the immediate mathematical content: \textbf{the proof difficulty in StdQM is not a feature of the underlying physics but a symptom of the formulation.}

Working in $3N$-dimensional configuration space with overlapping orbitals forces the singularities of the Coulomb potential to be controlled globally, requiring the elaborate apparatus of Dyson--Lenard--Lieb--Thirring. Working in real 3D space with non-overlapping unit densities localises the control, and stability falls out of the variational structure.

\paragraph{Pedagogical consequence.}
The natural pedagogical conclusion is that stability of matter, currently absent from standard quantum-mechanics curricula because the StdQM proof is too difficult, can be presented straightforwardly in the RealQM formulation as an \textbf{introductory result} --- a property of any system of non-overlapping unit electron densities under Coulomb interaction, derivable in a single page from the Hardy inequality.

This is not a small thing. Stability of matter is the precondition for the existence of bulk matter; it underlies the entire framework of chemistry, materials science, and biology. That it has been technically inaccessible to standard QM curricula is a structural problem, not a pedagogical one. The multiphase formulation removes the structural obstacle. A first-year student of RealQM can be shown why matter is stable.

\paragraph{What this is and is not.}
The RealQM proof of stability is not a refutation or correction of the StdQM proof. It is a proof of the analogous statement within a different mathematical framework. Both are correct within their domains; the RealQM proof is simpler because the framework is simpler. The fact that the simpler framework can carry the proof is evidence that the formulation is the right tool for understanding the physical phenomenon --- and the fact that this requires giving up antisymmetry in favour of non-overlap is precisely the architectural choice this book is about.

The Level-0 application of the framework --- protons and electrons as the constituent unit densities of the atomic nucleus, with the He-4 ground state and D + D $\to$ He-4 fusion as the worked examples --- was presented in Chapter~\ref{ch:nucleus} as part of the validation suite. The next chapter completes Part~IV with an honest accounting of the framework's regime across all levels of the hierarchy.

% --- End of Chapter 13 ---
